%% =====================================================================
%%  Capítulo 4 -- Sistemas bidimensionales: el plano de fase
%%  Formato nuevo (noviembre de 2026). Integra: Handout3.tex (versión anterior:
%%  plano de fase, diagonalización, teorema fundamental, linealización, los tres
%%  casos de Jordan, Duffing, equilibrios aislados, Hartman--Grobman y variedad
%%  central), las diapositivas SisBidim (2026), el cuaderno Planos_de_fase_Python
%%  de la clase del 19 de agosto y los scripts EjemploIntroSisBidim.m y
%%  CodigosVariosHandout3.m, reescritos en codigo/cap04_bidimensionales.
%%  Se agregaron: nulclinas, la clasificación completa traza--determinante,
%%  el diodo túnel de segundo orden y las variedades estable e inestable
%%  (clase del 31 de agosto). Correcciones: ver \S\ref{sec:bi-cierre}.
%%  Números verificados con herramientas/cap04_numeros.py (93/93) y
%%  codigo/cap04_bidimensionales/verificarCap04.m (38/38 en Octave 10.3).
%% =====================================================================

\capitulo[cap:bidimensionales]{Sistemas bidimensionales: el plano de fase}
  {Porque se sabe: siempre es más sencillo precisar dónde se está, pero la cosa se complica cuando se trata de detallar lo que se es.}
  {Rodrigo Fresán}

\begin{objetivos}[partes/04-bidimensionales.html]
  \item Dibujar e interpretar el plano de fase de un sistema de segundo orden: campo vectorial, trayectorias, nulclinas y equilibrios, sin resolver la ecuación.
  \item Resolver $\dot x=Ax$ con la exponencial de matriz y con un cambio de coordenadas, y probar que la solución es única.
  \item Clasificar el origen de $\dot x=Ax$ en $\R^2$ (nodo, silla, foco, centro y los casos degenerados) con la traza y el determinante, y saber cuáles de esos retratos sobreviven a una perturbación.
  \item Linealizar en un equilibrio y usar el teorema de Hartman--Grobman con sus hipótesis: qué garantiza, qué no garantiza y cuándo no dice nada.
  \item Reconocer las variedades estable e inestable de una silla y usarlas como fronteras entre cuencas de atracción.
  \item Reducir un sistema a su variedad central para decidir la estabilidad cuando la linealización tiene valores propios en el eje imaginario.
\end{objetivos}

%% ---------------------------------------------------------------------
\section{¿Por qué un capítulo para dos dimensiones?}
\label{sec:bi-porque}
%% ---------------------------------------------------------------------

En el \cref{cap:unidimensionales} el estado vivía en una recta, y una recta deja pocas opciones: las soluciones de $\dot x=f(x)$ son monótonas, van hacia un equilibrio o hacia infinito, y no pueden oscilar. Basta agregar una dimensión para que todo eso cambie. En el plano las trayectorias pueden girar alrededor de un equilibrio, acercarse por una dirección y alejarse por otra, o cerrarse sobre sí mismas. Los sistemas de este capítulo son
\begin{equation}
  \dot x_1=f_1(x_1,x_2),\qquad \dot x_2=f_2(x_1,x_2),\qquad\text{es decir,}\qquad \dot x=f(x),\quad x\in\R^2 ,
  \label{eq:bi-sistema}
\end{equation}
\emph{invariantes en el tiempo} ($f$ no depende de $t$) y sin entradas, con $f$ localmente Lipschitz, para que por cada condición inicial pase una sola solución (\cref{thm:pm-local}). Escribimos $\flujo t{x_0}$ para la solución que en $t=0$ pasa por $x_0$.

La pregunta del capítulo es local: \alerta{¿qué hacen las trayectorias cerca de un equilibrio?} La estrategia es la de siempre en ingeniería, linealizar, pero con la pregunta que casi nunca se hace en los cursos de control lineal: ¿cuándo dice la linealización la verdad sobre el sistema no lineal? La respuesta es el teorema de Hartman--Grobman, y sus excepciones (los equilibrios no hiperbólicos) son tan importantes como el teorema. Para ellas el capítulo cierra con la variedad central, que reduce el problema a una ecuación escalar como las del \cref{cap:unidimensionales}.

\notamargen{\textbf{Ruta de las clases.} Semana~2 (12 de agosto): \S\ref{sec:bi-plano}--\ref{sec:bi-clasificacion}. Semana~5 (31 de agosto), después de los fundamentos matemáticos: \S\ref{sec:bi-linealizacion}--\ref{sec:bi-central}. La prueba del teorema fundamental (\S\ref{sec:bi-lineales}) se lee por cuenta propia.}

Lo que pasa \emph{lejos} de los equilibrios (órbitas cerradas, ciclos límite, Poincaré--Bendixson) es el tema del \cref{cap:orbitas}, y la estabilidad sin linealizar, con funciones de Lyapunov, el del \cref{cap:lyapunov}. La \cref{fig:bi-mapa} muestra cómo encajan las piezas.

\begin{figure}[ht]
\centering
\includegraphics[width=0.95\linewidth]{cap04/t_mapa_bidim}
\caption{Mapa del capítulo. La linealización decide en los equilibrios hiperbólicos; en los demás, deciden los términos que la linealización desprecia.}
\label{fig:bi-mapa}
\end{figure}

\begin{historia}[La geometría antes que las fórmulas]
Henri Poincaré, en su memoria de 1881 sobre las curvas definidas por ecuaciones diferenciales~\citep{Poi1881}, cambió la pregunta: en lugar de buscar la fórmula de la solución, que casi nunca existe, estudió la \emph{forma} del conjunto de todas las soluciones. Los nombres de este capítulo son suyos: \emph{nœud} (nodo), \emph{col} (silla), \emph{foyer} (foco) y \emph{centre}. Que esa forma cerca de un equilibrio hiperbólico es la de la linealización lo probaron, de manera independiente, David Grobman en Moscú (1959) y Philip Hartman en Baltimore (1960)~\citep{Gro59,Har60}. La variedad central, para los casos que ellos dejan por fuera, viene de Pliss (1964) y Kelley (1967)~\citep{Pli64,Kel67}, y Carr la convirtió en una herramienta de cálculo~\citep{Car81}.
\end{historia}

\begin{pausa}
Antes de empezar, apueste: dos trayectorias distintas de \eqref{eq:bi-sistema}, ¿se pueden cruzar en el plano? ¿Y una trayectoria consigo misma?
\respuesta{Dos trayectorias distintas no se pueden cruzar ni tocar: en el punto común pasarían dos soluciones con la misma condición inicial, contra la unicidad. Una trayectoria sí puede volver a pasar por el mismo punto, pero entonces repite exactamente su recorrido: es una órbita cerrada (una solución periódica). Pueden, eso sí, \emph{acercarse} tanto como se quiera, por ejemplo las dos que llegan a un mismo equilibrio.}
\end{pausa}

%% ---------------------------------------------------------------------
\section{El plano de fase}
\label{sec:bi-plano}
%% ---------------------------------------------------------------------

La ecuación \eqref{eq:bi-sistema} asigna a cada punto $x$ del plano el vector $f(x)$: es un \clave{campo vectorial}. Una solución es una curva que en cada punto es tangente al campo, y se recorre con la velocidad que el campo indica.

\begin{definicion}[label=def:bi-plano]{Trayectoria, retrato y plano de fase}
La curva $\{\flujo t{x_0}: t\ge0\}$ (o $t\in\R$, si la solución existe para todo tiempo) es la \clave{trayectoria} u \clave{órbita} de \eqref{eq:bi-sistema} que pasa por $x_0$. El plano $x_1$--$x_2$ donde se dibujan es el \clave{plano de fase} (o plano de estado), y un dibujo con el campo y un conjunto representativo de trayectorias es un \clave{retrato de fase}.
\end{definicion}

El retrato se arma a mano en dos pasos. Primero se dibuja $f(x)$ en una malla de puntos: es la operación de \cmd{quiver} en \matlab. Por ejemplo, para
\begin{equation}
  \dot x_1=2x_1x_2,\qquad \dot x_2=x_1x_2 ,
  \label{eq:bi-ejcampo}
\end{equation}
en el punto $(1,2)$ el campo vale $f(1,2)=(4,2)$: la flecha sale de $(1,2)$ y apunta hacia $(5,4)$. Después se trazan trayectorias tangentes a las flechas. Si las flechas se dibujan con su largo real, $\norm{f(x)}$, las de las zonas rápidas tapan todo; casi siempre conviene \emph{normalizarlas} y leer solo la dirección.

\begin{figure}[ht]
\centering
\includegraphics[width=\linewidth]{cap04/f_campo}
\caption{Izquierda: el campo \eqref{eq:bi-ejcampo}, con flechas normalizadas. Las trayectorias son rectas paralelas a $(2,1)$ y los dos ejes están hechos de equilibrios. Derecha: la masa--resorte \eqref{eq:bi-resorte}: elipses alrededor del origen.}
\label{fig:bi-campo}
\end{figure}

\codigorepo{cap04_bidimensionales/C4_Ej1_CampoVectorial.m}{campo}{El campo vectorial con \cmd{quiver} y algunas trayectorias}

\begin{pausa}
En \eqref{eq:bi-ejcampo}, ¿por qué las trayectorias son rectas? ¿Hacia dónde se mueve un punto de la recta que pasa por $(1,2)$, y hasta dónde llega?
\respuesta{Porque $\frac{\dif}{\dif t}(x_1-2x_2)=2x_1x_2-2x_1x_2=0$: cada solución se queda en una recta $x_1-2x_2=c$. Sobre la recta $x_1-2x_2=-3$ que pasa por $(1,2)$, el producto $x_1x_2$ es positivo para $x_1>0$, así que el punto se aleja en la dirección $(2,1)$ (y escapa en tiempo finito, porque la velocidad crece cuadráticamente). Hacia atrás en el tiempo se acerca al equilibrio $(0,\tfrac32)$ sobre el eje $x_2$. Todos los puntos de los ejes son equilibrios: \emph{no} están aislados (\cref{sec:bi-linealizacion}).}
\end{pausa}

\subsection{Un ejemplo con física: la masa y el resorte}

Una masa $m$ colgada de un resorte lineal de constante $k$, sin fricción, cumple $m\ddot x+kx=0$. Con $x_1=x$ (posición) y $x_2=\dot x$ (velocidad),
\begin{equation}
  \dot x_1=x_2,\qquad \dot x_2=-\omega^2x_1,\qquad \omega^2=k/m .
  \label{eq:bi-resorte}
\end{equation}
Antes de resolver, se lee el campo por partes. Sobre el eje $x_1$ ($x_2=0$) el campo es $(0,-\omega^2x_1)$: vertical, hacia abajo si $x_1>0$ y hacia arriba si $x_1<0$, más largo cuanto más lejos del origen. Sobre el eje $x_2$ ($x_1=0$) es $(x_2,0)$: horizontal, hacia la derecha arriba y hacia la izquierda abajo. Juntando las dos cosas, el flujo gira \emph{en el sentido de las agujas del reloj}. Para saber si la vuelta se cierra basta una cuenta:
\[
  \frac{\dif}{\dif t}\bigl(\omega^2x_1^2+x_2^2\bigr)=2\omega^2x_1x_2+2x_2(-\omega^2x_1)=0 .
\]
Las trayectorias son las elipses $\omega^2x_1^2+x_2^2=c$: \clave{órbitas cerradas}, una para cada nivel de energía, y el origen es un equilibrio rodeado por ellas (\cref{fig:bi-campo}). Físicamente, sin fricción la energía $\tfrac12kx^2+\tfrac12m\dot x^2$ se conserva y la masa oscila para siempre con una amplitud que fija la condición inicial.

\subsection{Nulclinas: dónde el campo es vertical u horizontal}

Cuando la cuenta anterior no es posible, el retrato se organiza con las curvas donde una de las componentes del campo se anula.

\begin{definicion}[label=def:bi-nulclina]{Nulclinas}
La \clave{nulclina} de $x_1$ es el conjunto $\{x: f_1(x)=0\}$; sobre ella el campo es vertical. La de $x_2$ es $\{x: f_2(x)=0\}$; sobre ella el campo es horizontal. Los equilibrios son exactamente las intersecciones de una nulclina de $x_1$ con una de $x_2$.
\end{definicion}

Las nulclinas parten el plano en regiones donde los signos de $\dot x_1$ y $\dot x_2$ no cambian: en cada región el flujo va hacia uno de los cuatro cuadrantes (noreste, noroeste, \ldots), y eso ya da la forma gruesa del retrato.

\begin{ejemplo}[label=ej:bi-nulclinas]{El sistema de la clase del 19 de agosto}
Dibuje las nulclinas de
\begin{equation}
  \dot x_1=-x_1+x_1x_2,\qquad \dot x_2=x_1+x_2-2x_1x_2 ,
  \label{eq:bi-cuaderno}
\end{equation}
halle los equilibrios y describa el flujo en cada región.
\solucion $f_1=x_1(x_2-1)=0$ da dos rectas: $x_1=0$ y $x_2=1$. $f_2=0$ da $x_2(1-2x_1)=-x_1$, es decir, la hipérbola $x_2=x_1/(2x_1-1)$, con asíntotas $x_1=\tfrac12$ y $x_2=\tfrac12$. Las intersecciones son $(0,0)$ (con $x_1=0$) y $(1,1)$ (con $x_2=1$). En la región $x_1>1$, $x_2>1$ se tiene $\dot x_1>0$ y $\dot x_2<0$: el flujo va hacia el sureste. Rodeando $(1,1)$ los signos rotan, y el flujo gira alrededor de ese equilibrio; la \cref{sec:bi-linealizacion} dirá si se acerca en espiral o se aleja. El eje $x_2$ es invariante (si $x_1=0$, $\dot x_1=0$) y sobre él $\dot x_2=x_2$: el flujo se aleja del origen. Cerca del origen las trayectorias llegan en la dirección $(2,-1)$ y se van a lo largo del eje $x_2$ (\cref{fig:bi-nulclinas}); la \cref{sec:bi-linealizacion} explica de dónde sale esa dirección.
\end{ejemplo}

\begin{figure}[ht]
\centering
\includegraphics[width=0.78\linewidth]{cap04/f_nulclinas}
\caption{Nulclinas de \eqref{eq:bi-cuaderno}: de $x_1$ en coral (rectas $x_1=0$ y $x_2=1$), de $x_2$ en ámbar (la hipérbola). Los equilibrios están en las intersecciones de colores distintos.}
\label{fig:bi-nulclinas}
\end{figure}

\codigorepo{cap04_bidimensionales/C4_Ej2_Nulclinas.m}{nulclinas}{Nulclinas, equilibrios y retrato con la función \cmd{retratoFase} del curso}

\begin{trampa}[label=tr:bi-nulclinas]{``Dos nulclinas que se cruzan marcan un equilibrio''}
\plausible{En un equilibrio el campo es cero, y las nulclinas son justamente donde el campo se anula.}
Cada nulclina anula \emph{una} componente. El cruce de dos ramas de la \emph{misma} nulclina (en el ejemplo, las rectas $x_1=0$ y $x_2=1$ se cortan en $(0,1)$) no es un equilibrio: allí $f_1=0$, pero $f_2(0,1)=1\neq0$.
\correcto{Un equilibrio es un cruce de la nulclina de $x_1$ con la de $x_2$. Use colores distintos para las dos al dibujarlas.}
\end{trampa}

\subsection{Lo que el plano de fase no muestra}

El retrato elimina el tiempo: muestra \emph{por dónde} pasan las soluciones, no \emph{cuándo}. Dos sistemas con el mismo retrato pueden recorrerlo a velocidades muy distintas; si se multiplica $f$ por una función positiva, las trayectorias no cambian y los tiempos sí. Para el análisis cualitativo es justo lo que se quiere: rara vez importa cuánto tarda el sistema en llegar a un equilibrio, sino si llega y desde dónde.

\begin{trampa}[label=tr:bi-llega]{``La trayectoria llega al equilibrio''}
\plausible{En el retrato de un nodo estable las curvas terminan en el origen.}
Si $f$ es localmente Lipschitz, ninguna solución que empieza fuera de un equilibrio lo alcanza en tiempo finito: si $\flujo T{x_0}=\eq$, la solución que se queda en $\eq$ y la que llega a él pasarían por el mismo punto en $t=T$, contra la unicidad (hacia atrás en el tiempo). En $\dot x=-x$ la distancia se divide por $e$ cada unidad de tiempo, y nunca se anula. Sin Lipschitz sí puede pasar: $\dot x=-\sqrt{\abs x}$ llega a $0$ desde $x_0=1$ en $t=2$.
\correcto{Las trayectorias \emph{tienden} al equilibrio cuando $t\to\infty$. Las curvas del retrato que ``terminan'' en el equilibrio lo hacen en el límite.}
\end{trampa}

%% ---------------------------------------------------------------------
\section{Sistemas lineales: la exponencial de una matriz}
\label{sec:bi-lineales}
%% ---------------------------------------------------------------------

Cerca de un equilibrio, un sistema no lineal se parece (a veces) a uno lineal, así que primero hay que entender por completo
\begin{equation}
  \dot x=Ax,\qquad x(0)=x_0,\qquad A\in\R^{n\times n} .
  \label{eq:bi-lineal}
\end{equation}
La sección vale para cualquier $n$; la siguiente se especializa en $n=2$.

\subsection{La exponencial y el teorema fundamental}

Para una matriz cuadrada se define, como para los números,
\begin{equation}
  e^{At}=\sum_{k=0}^{\infty}\frac{A^kt^k}{k!}=I+At+\frac{A^2t^2}{2!}+\cdots
  \label{eq:bi-exponencial}
\end{equation}
La serie converge para toda $t$ porque cada término está acotado en norma por $\norm A^k\abs t^k/k!$, y $\sum\norm A^k\abs t^k/k!=e^{\norm A\abs t}<\infty$ (\cref{def:pm-inducida}). Más aún, converge \emph{uniformemente} en cada intervalo $\abs t\le T$, y lo mismo vale para la serie de las derivadas término a término.

\begin{lema}[label=lem:bi-derivada]{Derivada de la exponencial}
Para toda $t\in\R$, $\dfrac{\dif}{\dif t}e^{At}=Ae^{At}=e^{At}A$.
\end{lema}

\begin{demostracion}
Derivando término a término, lo que se permite porque la serie de las derivadas converge uniformemente en $\abs t\le T$,
\[
  \frac{\dif}{\dif t}\sum_{k=0}^{\infty}\frac{A^kt^k}{k!}=\sum_{k=1}^{\infty}\frac{A^kt^{k-1}}{(k-1)!}=A\sum_{j=0}^{\infty}\frac{A^jt^j}{j!}=Ae^{At} .
\]
Como $A$ conmuta con cada potencia $A^j$, también se puede sacar por la derecha: $Ae^{At}=e^{At}A$.
\end{demostracion}

\begin{teorema}[label=thm:bi-fundamental]{Teorema fundamental de los sistemas lineales}
Para cada $x_0\in\R^n$, el problema \eqref{eq:bi-lineal} tiene una única solución, definida para todo $t\in\R$:
\[
  x(t)=e^{At}x_0 .
\]
\end{teorema}

\begin{demostracion}
\emph{Es solución:} por el \cref{lem:bi-derivada}, $\dot x=Ae^{At}x_0=Ax(t)$, y $x(0)=Ix_0=x_0$.
\emph{Es la única:} sea $x(t)$ cualquier solución y $y(t)=e^{-At}x(t)$. Entonces
\[
  \dot y=-Ae^{-At}x+e^{-At}\dot x=-e^{-At}Ax+e^{-At}Ax=0 ,
\]
donde se usó que $e^{-At}$ conmuta con $A$. Luego $y$ es constante, $y(t)=y(0)=x_0$, y como $e^{At}e^{-At}=I$ (las dos conmutan), $x(t)=e^{At}y(t)=e^{At}x_0$.
\end{demostracion}

\notamargen{La unicidad también sale del \cref{thm:pm-local}, porque $f(x)=Ax$ es globalmente Lipschitz con constante $\norm A$; la prueba directa no necesita ese teorema.}

\begin{pausa}
¿Es cierto que $e^{At}$ es la matriz de las exponenciales $e^{a_{ij}t}$? ¿Y que $e^{(A+B)t}=e^{At}e^{Bt}$?
\respuesta{Lo primero solo si $A$ es diagonal: para $A=\bigl[\begin{smallmatrix}0&1\\0&0\end{smallmatrix}\bigr]$, $e^{At}=I+At=\bigl[\begin{smallmatrix}1&t\\0&1\end{smallmatrix}\bigr]$, y no $\bigl[\begin{smallmatrix}1&e^t\\1&1\end{smallmatrix}\bigr]$. Lo segundo solo si $AB=BA$; en general es falso, y en la prueba anterior hizo falta que $A$ conmute consigo misma.}
\end{pausa}

\subsection{Diagonalizar: cambiar a las coordenadas propias}

Calcular la serie no es práctico. El atajo es un cambio de coordenadas $x=Pz$ con $P$ invertible: si $\dot x=Ax$, entonces
\[
  \dot z=P^{-1}\dot x=P^{-1}APz ,
\]
y se escoge $P$ para que $P^{-1}AP$ sea lo más simple posible.

\begin{teorema}[label=thm:bi-diagonal]{Diagonalización \citep[\S1.2]{Per01}}
Si los valores propios $\lambda_1,\dots,\lambda_n$ de $A\in\R^{n\times n}$ son reales y distintos, cualquier conjunto de vectores propios correspondientes $v_1,\dots,v_n$ es una base de $\R^n$, la matriz $P=[v_1\ \cdots\ v_n]$ es invertible y
\[
  P^{-1}AP=\diag(\lambda_1,\dots,\lambda_n) .
\]
En consecuencia, $e^{At}=P\,\diag(e^{\lambda_1t},\dots,e^{\lambda_nt})\,P^{-1}$.
\end{teorema}

En las coordenadas $z$ el sistema se \emph{desacopla}: $\dot z_i=\lambda_iz_i$, $z_i(t)=z_i(0)e^{\lambda_it}$. Cada coordenada propia tiene su propio valor propio, y la solución en $x$ es una mezcla de esas exponenciales.

\begin{ejemplo}[label=ej:bi-diagonal]{Una matriz triangular}
Resuelva $\dot x_1=-x_1-3x_2$, $\dot x_2=2x_2$ por diagonalización.
\solucion $A=\bigl[\begin{smallmatrix}-1&-3\\0&2\end{smallmatrix}\bigr]$ es triangular, así que sus valores propios son $-1$ y $2$. Para $\lambda_1=-1$, $(A+I)v=0$ da $v_1=(1,0)$; para $\lambda_2=2$, $(A-2I)v=0$ da $-3v_{(1)}-3v_{(2)}=0$, $v_2=(-1,1)$. Con $P=\bigl[\begin{smallmatrix}1&-1\\0&1\end{smallmatrix}\bigr]$, $P^{-1}=\bigl[\begin{smallmatrix}1&1\\0&1\end{smallmatrix}\bigr]$ y $P^{-1}AP=\diag(-1,2)$. En coordenadas propias, $z_1(t)=z_1(0)e^{-t}$ y $z_2(t)=z_2(0)e^{2t}$. Volviendo a $x$,
\[
  e^{At}=\begin{bmatrix}e^{-t}&e^{-t}-e^{2t}\\0&e^{2t}\end{bmatrix},\qquad
  x_1(t)=x_1(0)e^{-t}+x_2(0)\bigl(e^{-t}-e^{2t}\bigr),\qquad x_2(t)=x_2(0)e^{2t}.
\]
(Compruebe: $\dot x_1+x_1=-3x_2(0)e^{2t}=-3x_2$.)
\end{ejemplo}

\codigorepo{cap04_bidimensionales/C4_Ej3_Exponencial.m}{diagonal}{\cmd{eig}, el cambio de coordenadas y la comparación con \cmd{expm}}

\begin{trampa}[label=tr:bi-diagonal]{``$\dot x_1=-x_1-3x_2$: como el coeficiente de $x_1$ es negativo, $x_1$ decae''}
\plausible{Si se ignora el acople, $\dot x_1=-x_1$ y $x_1(t)=x_1(0)e^{-t}$.}
El acople no se puede ignorar: $x_1(t)$ contiene el término $-x_2(0)e^{2t}$, que crece. Desde $x_0=(1,1)$, $x_1(t)=2e^{-t}-e^{2t}$ se anula en $t=\tfrac13\ln2\approx0.23$ y después se va a $-\infty$. Los elementos de la diagonal de $A$ no son sus valores propios (salvo si $A$ es triangular, como aquí, y aun así no dicen qué valor propio aparece en qué estado).
\correcto{La estabilidad de $\dot x=Ax$ la deciden los valores propios de $A$, y cada estado es, en general, una combinación de \emph{todos} los modos $e^{\lambda_it}$.}
\end{trampa}

\subsection{Subespacios estable, inestable y central}

Los vectores propios dicen además por dónde entran y salen las trayectorias. Si $v$ es un vector propio de $\lambda$ real, la recta $\{cv\}$ es \emph{invariante}: la solución que empieza en $x_0=cv$ es $ce^{\lambda t}v$, y se queda en la recta. Para valores propios complejos o repetidos se usan, respectivamente, las partes real e imaginaria de los vectores propios y los vectores propios generalizados, y se agrupan así.

\begin{definicion}[label=def:bi-subespacios]{Subespacios estable, inestable y central \citep[\S1.9]{Per01}}
Sea $\{w_1,\dots,w_n\}$ una base de $\R^n$ formada por vectores propios generalizados de $A$ (para cada par complejo, sus partes real e imaginaria). Los subespacios
\[
  E^s=\operatorname{gen}\{w_j:\Re\lambda_j<0\},\qquad E^u=\operatorname{gen}\{w_j:\Re\lambda_j>0\},\qquad E^c=\operatorname{gen}\{w_j:\Re\lambda_j=0\}
\]
son el \clave{subespacio estable}, el \clave{inestable} y el \clave{central}. Son invariantes bajo $e^{At}$ y $\R^n=E^s\oplus E^u\oplus E^c$.
\end{definicion}

Las soluciones que empiezan en $E^s$ tienden al origen, las de $E^u$ tienden al origen hacia atrás en el tiempo, y las de $E^c$ ni crecen ni decaen exponencialmente. En el \cref{ej:bi-diagonal}, $E^s$ es el eje $x_1$ y $E^u$ la recta generada por $(-1,1)$: solo las condiciones iniciales sobre el eje $x_1$ ($x_2(0)=0$) llevan al origen. La \cref{sec:bi-variedades} muestra qué queda de $E^s$ y $E^u$ cuando el sistema no es lineal.

%% ---------------------------------------------------------------------
\section{Los retratos de fase lineales en el plano}
\label{sec:bi-clasificacion}
%% ---------------------------------------------------------------------

En $\R^2$ toda matriz real es semejante a una de tres \clave{formas de Jordan reales}: existe $P$ invertible con $P^{-1}AP=J$ y
\begin{equation}
  J_1=\begin{bmatrix}\lambda_1&0\\0&\lambda_2\end{bmatrix},\qquad
  J_2=\begin{bmatrix}\lambda&1\\0&\lambda\end{bmatrix},\qquad
  J_3=\begin{bmatrix}\alpha&-\beta\\\beta&\alpha\end{bmatrix}\ (\beta>0) .
  \label{eq:bi-jordan}
\end{equation}
$J_1$ corresponde a valores propios reales con dos vectores propios independientes (distintos, o iguales con $A=\lambda I$); $J_2$, a un valor propio real doble con un solo vector propio; $J_3$, al par complejo $\alpha\pm j\beta$. Como $x=Pz$ es una deformación lineal del plano (estira, rota, refleja), el retrato en $x$ es el retrato en $z$ visto a través de $P$: rectas van a rectas y el origen al origen. Por eso basta estudiar las tres formas. Supondremos $\det A\neq0$ hasta el final de la sección.

\subsection{Valores propios reales distintos: nodos y sillas}

Con $J_1$, $z_1(t)=z_1(0)e^{\lambda_1t}$ y $z_2(t)=z_2(0)e^{\lambda_2t}$. Eliminando $t$ (si $z_1(0)\neq0$),
\begin{equation}
  z_2=c\,\abs{z_1}^{\lambda_2/\lambda_1},\qquad c=\frac{z_2(0)}{\abs{z_1(0)}^{\lambda_2/\lambda_1}} .
  \label{eq:bi-nodo}
\end{equation}
\begin{itemize}
  \item $\lambda_1,\lambda_2<0$: \clave{nodo estable}. El exponente $\lambda_2/\lambda_1$ es positivo y las curvas son ``parábolas'' que llegan al origen. Si $\abs{\lambda_1}>\abs{\lambda_2}$ ($z_1$ es la coordenada rápida), $\lambda_2/\lambda_1<1$ y las curvas llegan \emph{tangentes al eje $z_2$}: primero desaparece la componente rápida y la trayectoria termina a lo largo de la dirección \emph{lenta}.
  \item $\lambda_1,\lambda_2>0$: \clave{nodo inestable}. Las mismas curvas, recorridas hacia afuera; salen tangentes a la dirección lenta (la de menor $\lambda$).
  \item $\lambda_1<0<\lambda_2$: \clave{silla}. El exponente es negativo y las curvas son ``hipérbolas''. Solo las condiciones iniciales sobre $E^s$ (el eje $z_1$) llegan al origen; todas las demás terminan escapando a lo largo de $E^u$.
\end{itemize}
Si $\lambda_1=\lambda_2=\lambda$ y $A=\lambda I$, todas las rectas por el origen son invariantes: el retrato es una \clave{estrella} (nodo propio).

\begin{figure}[ht]
\centering
\includegraphics[width=\linewidth]{cap04/f_retratos}
\caption{Los retratos lineales en las coordenadas propias. Arriba: nodo estable ($\lambda=-6,-2$), silla ($-1,2$), estrella ($-1,-1$). Abajo: nodo impropio ($J_2$, $\lambda=-1$), foco estable ($-0.5\pm2j$), centro ($\pm2j$). En azul, los subespacios propios.}
\label{fig:bi-retratos}
\end{figure}

\codigorepo{cap04_bidimensionales/C4_Ej4_Galeria.m}{galeria}{Los seis retratos de la \cref{fig:bi-retratos}}

\begin{trampa}[label=tr:bi-rapida]{``En un nodo, las trayectorias llegan por la dirección rápida''}
\plausible{La dirección rápida es la que más atrae, así que debería dominar la llegada.}
Es al revés. Con $\lambda_1=-6$ y $\lambda_2=-2$, la componente $z_1$ decae como $e^{-6t}$ y la $z_2$ como $e^{-2t}$: cuando la trayectoria está cerca del origen, lo que queda es $z_2$. En \eqref{eq:bi-nodo} el exponente es $\tfrac13$ y $z_2\propto\abs{z_1}^{1/3}$, una curva tangente al eje $z_2$ en el origen.
\correcto{Las trayectorias de un nodo estable llegan tangentes al vector propio \emph{lento} (el de $\abs\lambda$ menor), salvo las que empiezan sobre el vector propio rápido.}
\end{trampa}

\subsection{Un valor propio doble con un solo vector propio}

Con $J_2$ (es decir, $A\neq\lambda I$ con $\lambda$ doble), $\dot z_1=\lambda z_1+z_2$, $\dot z_2=\lambda z_2$, de modo que
\[
  z_2(t)=z_2(0)e^{\lambda t},\qquad z_1(t)=\bigl(z_1(0)+z_2(0)\,t\bigr)e^{\lambda t} .
\]
Eliminando $t=\frac1\lambda\ln\frac{z_2}{z_2(0)}$,
\[
  z_1=z_2\left(\frac{z_1(0)}{z_2(0)}+\frac1\lambda\ln\frac{z_2}{z_2(0)}\right) .
\]
Hay una sola recta invariante (el eje $z_1$, generado por el vector propio), y todas las trayectorias llegan tangentes a ella, después de dar media vuelta: es un \clave{nodo impropio} (o degenerado), estable si $\lambda<0$ e inestable si $\lambda>0$. Se lo puede ver como el caso límite entre un nodo y un foco.

\subsection{Valores propios complejos: focos y centros}

Si $\lambda=\alpha+j\beta$ ($\beta>0$) tiene vector propio $v=a+jb$, con $P=[a\ \ -b]$ se obtiene $J_3$ (compruébelo multiplicando $AP$). En coordenadas polares, $z_1=r\cos\theta$, $z_2=r\sin\theta$,
\[
  \dot z_1=\alpha z_1-\beta z_2,\quad \dot z_2=\beta z_1+\alpha z_2
  \quad\Longrightarrow\quad
  \dot r=\alpha r,\qquad \dot\theta=\beta .
\]
El ángulo crece a tasa constante y el radio cambia como $e^{\alpha t}$:
\begin{itemize}
  \item $\alpha<0$: \clave{foco estable} (espirales que entran).
  \item $\alpha>0$: \clave{foco inestable} (espirales que salen).
  \item $\alpha=0$: \clave{centro}. Ni entran ni salen: el origen queda rodeado por órbitas cerradas (círculos en $z$, elipses en $x$), como la masa--resorte.
\end{itemize}

\begin{trampa}[label=tr:bi-giro]{``El foco gira en sentido antihorario porque $\dot\theta=\beta>0$''}
\plausible{En las coordenadas de Jordan, $\dot\theta=\beta>0$.}
Eso vale en $z$. El cambio $x=Pz$ puede incluir una reflexión ($\det P<0$), y entonces invierte el sentido. La masa--resorte \eqref{eq:bi-resorte} tiene $\beta=\omega>0$ y gira en sentido horario en el plano $x$.
\correcto{El sentido de giro en el plano $x$ se lee directamente en $A$: en el punto $(1,0)$ la velocidad es $(a_{11},a_{21})$. Si $a_{21}>0$ el giro es antihorario; si $a_{21}<0$, horario. (Con valores propios complejos, $a_{21}\neq0$ siempre.)}
\end{trampa}

\subsection{Traza y determinante}

No hace falta calcular $J$: la traza $\tau=\tr A=\lambda_1+\lambda_2$ y el determinante $\delta=\det A=\lambda_1\lambda_2$ bastan, porque los valores propios son
\[
  \lambda=\frac{\tau\pm\sqrt{\tau^2-4\delta}}{2} .
\]

\begin{teorema}[label=thm:bi-trazadet]{Clasificación por traza y determinante \citep[\S1.5]{Per01}}
Sea $A\in\R^{2\times2}$ con $\delta=\det A\neq0$ y $\tau=\tr A$. El origen es el único equilibrio de $\dot x=Ax$, y es
\begin{enumerate}
  \item una \textbf{silla} si $\delta<0$;
  \item un \textbf{nodo} si $\delta>0$ y $\tau^2-4\delta\ge0$ (estrella o impropio si $\tau^2=4\delta$), estable si $\tau<0$ e inestable si $\tau>0$;
  \item un \textbf{foco} si $\delta>0$, $\tau^2-4\delta<0$ y $\tau\neq0$, estable si $\tau<0$ e inestable si $\tau>0$;
  \item un \textbf{centro} si $\delta>0$ y $\tau=0$.
\end{enumerate}
\end{teorema}

\begin{demostracion}
Como $\delta\neq0$, $Ax=0$ solo si $x=0$. (1) $\lambda_1\lambda_2=\delta<0$: dos valores propios reales de signo opuesto. (2) Reales (discriminante $\ge0$) con producto positivo: del mismo signo, el de su suma $\tau$; además $\tau^2\ge4\delta>0$ garantiza $\tau\neq0$. (3) Complejos con parte real $\alpha=\tau/2\neq0$. (4) Complejos con parte real nula, $\lambda=\pm j\sqrt\delta$.
\end{demostracion}

\begin{figure}[ht]
\centering
\includegraphics[width=0.85\linewidth]{cap04/f_traza_det}
\caption{El plano traza--determinante. La parábola $\tau^2=4\delta$ separa nodos y focos; el semieje $\tau=0$, $\delta>0$ son los centros. Los retratos de las fronteras (centros, estrellas, nodos impropios, $\delta=0$) no son robustos.}
\label{fig:bi-trazadet}
\end{figure}

\video{V4_1_traza_determinante}{Un punto recorre el plano traza--determinante}
\codigorepo{cap04_bidimensionales/C4_Ej5_TrazaDeterminante.m}{clasificar}{La función \cmd{clasificarEquilibrio} del curso}

\begin{ejemplo}[label=ej:bi-barrido]{Un parámetro que recorre todos los casos}
Para $\dot x_1=ax_1$, $\dot x_2=-x_2$, describa el retrato cuando $a$ va de $-\infty$ a $\infty$.
\solucion $x_1=x_1(0)e^{at}$, $x_2=x_2(0)e^{-t}$; $\tau=a-1$, $\delta=-a$. Si $a<-1$: nodo estable, con $x_2$ como dirección lenta (las curvas llegan tangentes al eje $x_2$). Si $a=-1$: estrella. Si $-1<a<0$: nodo estable con $x_1$ como dirección lenta (llegan tangentes al eje $x_1$). Si $a=0$: $\delta=0$, todo el eje $x_1$ está hecho de equilibrios y las trayectorias son rectas verticales que bajan hacia él. Si $a>0$: silla. El paso por $a=0$ es una bifurcación (\cref{fig:bi-barrido}).
\end{ejemplo}

\begin{figure}[ht]
\centering
\includegraphics[width=\linewidth]{cap04/f_barrido}
\caption{$\dot x_1=ax_1$, $\dot x_2=-x_2$ para $a=-4$, $-1$, $-0.3$, $0$ y $1$. En $a=-1$ la dirección lenta cambia de eje; en $a=0$ aparece una recta de equilibrios.}
\label{fig:bi-barrido}
\end{figure}

\subsection{Qué retratos sobreviven a una perturbación}

Los coeficientes de un modelo nunca se conocen exactamente. Si se cambia $A$ por $A+\Delta A$ con $\norm{\Delta A}$ pequeño, los valores propios se mueven poco (dependen continuamente de los elementos de la matriz~\citep{GoVa83}). Entonces:
\begin{itemize}
  \item un valor propio con parte real negativa (positiva) la conserva: \alerta{nodos, focos y sillas se conservan}. Son retratos \clave{estructuralmente estables};
  \item un par $\pm j\beta$ puede moverse a cualquiera de los dos semiplanos: \alerta{un centro se vuelve foco estable o inestable} con una perturbación arbitrariamente pequeña ($A+\varepsilon I$ tiene $\Re\lambda=\varepsilon$);
  \item un valor propio doble se puede partir en dos reales o en un par complejo: \alerta{un nodo impropio o una estrella se vuelve nodo o foco} (del mismo tipo de estabilidad). Por ejemplo, $\bigl[\begin{smallmatrix}-1&1\\-\varepsilon&-1\end{smallmatrix}\bigr]$ tiene $\lambda=-1\pm j\sqrt\varepsilon$;
  \item un valor propio cero se vuelve positivo o negativo, y la recta de equilibrios desaparece.
\end{itemize}

\begin{figure}[ht]
\centering
\includegraphics[width=\linewidth]{cap04/f_estructural}
\caption{Izquierda: un centro y dos perturbaciones de tamaño $0.05$. Derecha: un nodo impropio y la perturbación $\varepsilon=0.05$ que lo vuelve foco. La estabilidad se conserva en el segundo caso, no en el primero.}
\label{fig:bi-estructural}
\end{figure}

\codigorepo{cap04_bidimensionales/C4_Ej6_Estructural.m}{perturbar}{Perturbar un centro y un nodo impropio}

\begin{pausa}
¿Es cierto que todo sistema plano tiene al menos un equilibrio? Piense en $\dot x_1=\alpha+\sin(x_1+x_2)+x_1$, $\dot x_2=\alpha+\sin(x_1+x_2)-x_1$.
\respuesta{No. Restando las dos ecuaciones, un equilibrio necesita $x_1=0$, y entonces $\sin x_2=-\alpha$. Si $\alpha>1$ no hay solución: el sistema no tiene equilibrios (con $\alpha\le1$ sí los tiene, infinitos). Los sistemas lineales siempre tienen al origen; los no lineales, no necesariamente.}
\end{pausa}

%% ---------------------------------------------------------------------
\section{Linealización y el teorema de Hartman--Grobman}
\label{sec:bi-linealizacion}
%% ---------------------------------------------------------------------

Sea $\eq$ un equilibrio de \eqref{eq:bi-sistema}, $f(\eq)=0$, con $f$ de clase $\Cont1$. Con la variable de desviación $e=x-\eq$, el desarrollo de Taylor de primer orden (\cref{def:pm-jacobiano}) da
\begin{equation}
  \dot e=f(\eq+e)=\underbrace{f(\eq)}_{=0}+Ae+g(e),\qquad
  A=\jacat fx{x=\eq}=\begin{bmatrix}\partial f_1/\partial x_1&\partial f_1/\partial x_2\\[2pt]\partial f_2/\partial x_1&\partial f_2/\partial x_2\end{bmatrix}_{x=\eq},
  \label{eq:bi-linealizacion}
\end{equation}
donde el resto cumple $\norm{g(e)}/\norm e\to0$ cuando $e\to0$. La \clave{linealización} en $\eq$ es $\dot e=Ae$. La pregunta es cuándo el retrato de $\dot e=Ae$, que la \cref{sec:bi-clasificacion} sabe dibujar, describe al de $\dot x=f(x)$ cerca de $\eq$.

\subsection{Equilibrios aislados}

\begin{definicion}[label=def:bi-aislado]{Equilibrio aislado}
Un equilibrio $\eq$ es \clave{aislado} si existe $r>0$ tal que no hay otro equilibrio en la bola $\norm{x-\eq}<r$.
\end{definicion}

\begin{proposicion}[label=prop:bi-aislado]{Jacobiana invertible implica equilibrio aislado}
Si $f\in\Cont1$, $f(\eq)=0$ y $A=\jacat fx{\eq}$ es invertible, entonces $\eq$ es un equilibrio aislado.
\end{proposicion}

\begin{demostracion}
Como $A$ es invertible, $m=\min_{\norm v=1}\norm{Av}>0$ (el mínimo de una función continua y positiva sobre un compacto), y $\norm{Ae}\ge m\norm e$ para todo $e$. Por la definición de derivada, existe $r>0$ tal que $\norm{g(e)}\le\tfrac m2\norm e$ si $\norm e<r$. Entonces, para $0<\norm e<r$,
\[
  \norm{f(\eq+e)}=\norm{Ae+g(e)}\ge m\norm e-\tfrac m2\norm e=\tfrac m2\norm e>0 ,
\]
y $\eq+e$ no es un equilibrio.
\end{demostracion}

El recíproco es falso: $\dot x=x^3$ tiene un equilibrio aislado en $0$ y $f'(0)=0$. Y si la jacobiana es singular puede haber una curva entera de equilibrios, como los ejes de \eqref{eq:bi-ejcampo}. La hipótesis que importa para el retrato no es solo $\det A\neq0$ sino una más fuerte.

\begin{definicion}[label=def:bi-hiperbolico]{Equilibrio hiperbólico}
El equilibrio $\eq$ es \clave{hiperbólico} si ningún valor propio de $A=\jacat fx{\eq}$ tiene parte real nula. En $\R^2$: nodos, focos y sillas son hiperbólicos; centros y equilibrios con un valor propio cero no lo son.
\end{definicion}

\subsection{El teorema}

\begin{teorema}[label=thm:bi-hg]{Hartman--Grobman \citep{Gro59,Har60}, \citep[\S2.8]{Per01}}
Sea $f\in\Cont1(\D)$ y $\eq\in\D$ un equilibrio \emph{hiperbólico} de $\dot x=f(x)$, con $A=\jacat fx{\eq}$. Existen vecindades $U$ de $\eq$ y $V$ del origen y un homeomorfismo $h:U\to V$ (una biyección continua con inversa continua) tal que
\[
  h\bigl(\flujo t{x_0}\bigr)=e^{At}h(x_0)
\]
para todo $x_0\in U$ y todo $t$ con $\flujo t{x_0}\in U$. Es decir, $h$ lleva las trayectorias del sistema no lineal en trayectorias de la linealización, conservando el sentido del tiempo.
\end{teorema}

\begin{figure}[ht]
\centering
\includegraphics[width=\linewidth]{cap04/f_hartman}
\caption{Hartman--Grobman en la silla de Duffing (\S\ref{sec:bi-ejemplos}). Izquierda: cerca del equilibrio las trayectorias no lineales (azul) se parecen a las de la linealización (gris), con las separatrices curvadas. Derecha: lejos del equilibrio la semejanza se pierde; el teorema no dice nada allí.}
\label{fig:bi-hartman}
\end{figure}

\paragraph{Qué dice y qué no dice.}
\begin{enumerate}
  \item Es \emph{local}: no dice qué tan grande es $U$, ni nada de lo que pasa fuera de ella.
  \item $h$ es continua pero, en general, no diferenciable, y no es única. Se puede deformar el plano cerca del equilibrio, pero no se puede calcular $h$ en la práctica. Lo que da el teorema es una certeza cualitativa: las trayectorias no lineales tienen la misma \emph{topología} que las lineales.
  \item Una consecuencia es el método indirecto de Lyapunov (\cref{thm:ly-indirecto}): si $A$ es Hurwitz, $\eq$ es asintóticamente estable; si algún valor propio tiene parte real positiva, es inestable.
  \item Si $\eq$ \emph{no} es hiperbólico, el teorema no dice nada. No dice ``entonces el retrato es otro''; dice ``no sé''.
\end{enumerate}

\begin{trampa}[label=tr:bi-hgfoco]{``Hartman--Grobman dice si el equilibrio no lineal es un nodo o un foco''}
\plausible{El teorema dice que el retrato no lineal ``es'' el de la linealización.}
Lo es \emph{salvo homeomorfismo}, y un homeomorfismo puede enderezar espirales: un nodo estable y un foco estable son topológicamente equivalentes entre sí (los dos son ``todo llega al punto''). Hartman--Grobman, por sí solo, distingue atractor, repulsor y silla, no más. Que las trayectorias no lineales giren alrededor de un foco lineal (y no lleguen tangentes a una dirección en un nodo) es cierto, pero viene de un resultado más fino: en el plano, si $f\in\Cont2$, la conjugación puede tomarse de clase $\Cont1$~\citep{Har60b}, y un difeomorfismo sí conserva esa distinción.
\correcto{Para la estabilidad y el número de direcciones que entran y salen, basta Hartman--Grobman. Para decir ``nodo'' o ``foco'' del sistema no lineal, se necesita $f\in\Cont2$ y el resultado de Hartman para el plano. En la práctica (y en este curso) se usa el nombre de la linealización.}
\end{trampa}

\subsection{Ejemplos}

\begin{ejemplo}[label=ej:bi-cuadernoeq]{Los equilibrios del sistema de la clase}
Clasifique los equilibrios de \eqref{eq:bi-cuaderno}.
\solucion $A(x)=\bigl[\begin{smallmatrix}x_2-1&x_1\\1-2x_2&1-2x_1\end{smallmatrix}\bigr]$. En $(0,0)$, $A=\bigl[\begin{smallmatrix}-1&0\\1&1\end{smallmatrix}\bigr]$ con $\delta=-1<0$: silla, con $E^s$ generado por $(2,-1)$ y $E^u$ por $(0,1)$ (la dirección que se vio en el \cref{ej:bi-nulclinas}). En $(1,1)$, $A=\bigl[\begin{smallmatrix}0&1\\-1&-1\end{smallmatrix}\bigr]$ con $\tau=-1$, $\delta=1$ y $\tau^2-4\delta=-3<0$: foco estable, $\lambda=-\tfrac12\pm j\tfrac{\sqrt3}2$, y gira en sentido horario ($a_{21}=-1$). Ambos son hiperbólicos: el retrato no lineal tiene, cerca de cada uno, esa forma.
\end{ejemplo}

\begin{ejemplo}[label=ej:bi-tres]{Tres equilibrios (diapositivas 2026)}
Halle y clasifique los equilibrios de
\begin{equation}
  \dot x_1=x_1+x_1x_2,\qquad \dot x_2=-x_2+x_2^2+x_1x_2-x_1^3 .
  \label{eq:bi-tres}
\end{equation}
\solucion $f_1=x_1(1+x_2)=0$. Si $x_1=0$: $x_2(x_2-1)=0$, equilibrios $(0,0)$ y $(0,1)$. Si $x_2=-1$: $2-x_1-x_1^3=0$, cuya única raíz real es $x_1=1$; equilibrio $(1,-1)$. La jacobiana es $A(x)=\bigl[\begin{smallmatrix}1+x_2&x_1\\x_2-3x_1^2&-1+2x_2+x_1\end{smallmatrix}\bigr]$:
\begin{itemize}
  \item $(0,0)$: $\diag(1,-1)$, silla.
  \item $(0,1)$: $\bigl[\begin{smallmatrix}2&0\\1&1\end{smallmatrix}\bigr]$, $\lambda=2,1$: nodo inestable.
  \item $(1,-1)$: $\bigl[\begin{smallmatrix}0&1\\-4&-2\end{smallmatrix}\bigr]$, $\tau=-2$, $\delta=4$: foco estable, $\lambda=-1\pm j\sqrt3$.
\end{itemize}
Los tres son hiperbólicos (\cref{fig:bi-tres}).
\end{ejemplo}

\begin{figure}[ht]
\centering
\includegraphics[width=0.8\linewidth]{cap04/f_tres}
\caption{Retrato de \eqref{eq:bi-tres}. Cerca de cada equilibrio el retrato es el de su linealización; lejos, las trayectorias que salen del nodo $(0,1)$ se reparten entre el foco y el infinito.}
\label{fig:bi-tres}
\end{figure}

\codigorepo{cap04_bidimensionales/C4_Ej7_Linealizacion.m}{tres}{Equilibrios con \cmd{fsolve}, jacobianas y clasificación}

\begin{ejemplo}[label=ej:bi-epsilon]{Cuando la linealización es un centro}
Estudie el origen de $\dot x_1=x_2$, $\dot x_2=-x_1-\varepsilon x_1^2x_2$.
\solucion La jacobiana en el origen es $\bigl[\begin{smallmatrix}0&1\\-1&0\end{smallmatrix}\bigr]$, $\lambda=\pm j$, para \emph{cualquier} $\varepsilon$: un centro, no hiperbólico, y Hartman--Grobman calla. Con $V=x_1^2+x_2^2$,
\[
  \Vdot=2x_1x_2+2x_2(-x_1-\varepsilon x_1^2x_2)=-2\varepsilon x_1^2x_2^2 .
\]
Si $\varepsilon>0$, $V$ no crece: las trayectorias no se alejan, y de hecho se acercan en espiral al origen (desde $r=0.5$ llegan a $r\approx0.23$ en $t=60$, lentamente). Si $\varepsilon<0$, $V$ no decrece y las trayectorias se alejan. El término cúbico, que la linealización desprecia, decide (\cref{fig:bi-epsilon}). Para concluir estabilidad asintótica cuando $\varepsilon>0$ falta un argumento, porque $\Vdot$ se anula en los ejes: es el principio de invarianza de LaSalle (capítulo~8).
\end{ejemplo}

\begin{figure}[ht]
\centering
\includegraphics[width=\linewidth]{cap04/f_epsilon}
\caption{La misma linealización (un centro) y tres retratos no lineales distintos: $\varepsilon=1$ (espiral lenta hacia adentro), $\varepsilon=0$ (centro de verdad) y $\varepsilon=-1$ (espiral hacia afuera).}
\label{fig:bi-epsilon}
\end{figure}

\codigorepo{cap04_bidimensionales/C4_Ej8_Centro.m}{epsilon}{El centro lineal y lo que hace el término cúbico}

\begin{trampa}[label=tr:bi-centro]{``La linealización es un centro: el sistema no lineal oscila alrededor del equilibrio''}
\plausible{Cerca del equilibrio el sistema se parece a su linealización, y la linealización tiene órbitas cerradas.}
El parecido solo está garantizado si el equilibrio es hiperbólico, y un centro no lo es. En el \cref{ej:bi-epsilon} la misma linealización corresponde a un equilibrio asintóticamente estable, a un centro y a uno inestable, según el signo de $\varepsilon$. Lo mismo pasa en una dimensión: $\dot x=-x^3$ y $\dot x=x^3$ tienen la linealización $\dot x=0$ y se comportan al revés.
\correcto{Si la linealización tiene valores propios con parte real nula, el retrato no lineal lo deciden los términos de orden superior. Hay que usar una función de Lyapunov (\cref{cap:lyapunov}), la variedad central (\cref{sec:bi-central}) o una cantidad conservada (\cref{cap:orbitas}).}
\end{trampa}

%% ---------------------------------------------------------------------
\section{Dos sistemas con varios equilibrios}
\label{sec:bi-ejemplos}
%% ---------------------------------------------------------------------

Los dos ejemplos de esta sección tienen dos equilibrios estables separados por una silla. Esa configuración, la \clave{biestabilidad}, es la forma más simple de memoria que tiene un sistema físico: el estado final depende de dónde empezó.

\subsection{La ecuación de Duffing}

\begin{historia}[Georg Duffing (1861--1944)]
Duffing fue un ingeniero mecánico alemán, formado en Karlsruhe, que se dedicó por su cuenta al estudio de las vibraciones. En 1910 hizo una pasantía en la Westinghouse, en Estados Unidos, y en 1913 se estableció en Berlín como inventor independiente. En 1918 publicó \emph{Erzwungene Schwingungen bei veränderlicher Eigenfrequenz und ihre technische Bedeutung}~\citep{Duf18}, sobre oscilaciones forzadas con frecuencia natural variable, donde aparece la ecuación que lleva su nombre. Con dificultades económicas después de la guerra, en 1921 entró a dirigir el laboratorio de la Ölwerke Stern-Sonneborn (Ossag) en Hamburgo, donde trabajó en la viscosidad y la fricción de los lubricantes~\citep{KoBr11}. Su ecuación se volvió famosa medio siglo después, cuando la versión forzada se convirtió en uno de los ejemplos clásicos de caos.
\end{historia}

La forma general de la ecuación es
\[
  \ddot x+\delta\dot x+\alpha x+\beta x^3=\gamma\cos(\omega t),
\]
donde $\delta\ge0$ es el amortiguamiento, $\alpha$ la rigidez lineal, $\beta$ la no linealidad del resorte ($\beta>0$, resorte que se endurece; $\beta<0$, que se ablanda) y $\gamma\cos\omega t$ el forzamiento. Sin forzamiento y con $\alpha=-1$, $\beta=1$, describe una viga de acero entre dos imanes: un pozo de potencial doble. Con $x_1=x$, $x_2=\dot x$,
\begin{equation}
  \dot x_1=x_2,\qquad \dot x_2=x_1-x_1^3-\delta x_2 .
  \label{eq:bi-duffing}
\end{equation}

\begin{ejemplo}[label=ej:bi-duffing]{Duffing amortiguado}
Halle y clasifique los equilibrios de \eqref{eq:bi-duffing} para $\delta>0$.
\solucion $x_2=0$ y $x_1(1-x_1^2)=0$: equilibrios en $(0,0)$ y $(\pm1,0)$. La jacobiana es $\bigl[\begin{smallmatrix}0&1\\1-3x_1^2&-\delta\end{smallmatrix}\bigr]$.
\begin{itemize}
  \item $(0,0)$: $\delta_A=\det A=-1<0$, silla, con $\lambda=\tfrac12\bigl(-\delta\pm\sqrt{\delta^2+4}\bigr)$.
  \item $(\pm1,0)$: $\tau=-\delta<0$, $\det A=2>0$, $\tau^2-4\det A=\delta^2-8$. Son \emph{focos} estables si $\delta<2\sqrt2\approx2.83$ y \emph{nodos} estables si $\delta\ge2\sqrt2$.
\end{itemize}
Con $\delta=0.5$: silla con $\lambda=0.781$ y $-1.281$, focos con $\lambda=-0.25\pm1.392j$. Los tres son hiperbólicos.
\end{ejemplo}

La energía $E=\tfrac12x_2^2-\tfrac12x_1^2+\tfrac14x_1^4$ (cinética más potencial de doble pozo) cumple $\dot E=-\delta x_2^2\le0$: el amortiguamiento disipa energía, y casi todas las trayectorias terminan en el fondo de uno de los dos pozos. ¿En cuál? La respuesta la da la silla y se ve en la \cref{sec:bi-variedades}: desde $(1.5,0)$ la viga termina en $(1,0)$, pero desde $(2,0)$, que tiene más energía, cruza la barrera y termina en $(-1,0)$.

\codigorepo{cap04_bidimensionales/C4_Ej9_Duffing.m}{duffing}{Equilibrios de Duffing en función de $\delta$ y dos condiciones iniciales vecinas}

\begin{trampa}[label=tr:bi-caos]{``La ecuación de Duffing sin amortiguamiento ni forzamiento es caótica''}
\plausible{Duffing es el ejemplo clásico de caos, y sin fricción nada frena las oscilaciones.}
Con $\delta=\gamma=0$ el sistema es plano e invariante en el tiempo, y conserva $E$: cada trayectoria vive en una curva de nivel de $E$, y el retrato es un ocho rodeado de órbitas cerradas (\cref{ej:oc-duffing}). Las trayectorias de un sistema autónomo en el plano no pueden ser caóticas: el teorema de Poincaré--Bendixson (\cref{cap:orbitas}) solo deja equilibrios, órbitas cerradas y ciclos de equilibrios como comportamiento a largo plazo.
\correcto{El caos de Duffing aparece con el forzamiento $\gamma\cos\omega t$: el sistema forzado tiene un estado de dimensión tres ($x_1$, $x_2$ y la fase del forzamiento), y en tres dimensiones Poincaré--Bendixson ya no aplica.}
\end{trampa}

\subsection{El diodo túnel con su inductancia}

En la \cref{sec:un-diodo} se estudió el circuito del diodo túnel sin inductancia. Con la inductancia $L$ en serie con la fuente, el estado son el voltaje del condensador y la corriente del inductor (\cref{prob:in-diodo}), y el modelo de \citet[\S1.2.2]{Kha02} es
\begin{equation}
  \dot x_1=\frac1C\bigl(-h(x_1)+x_2\bigr),\qquad \dot x_2=\frac1L\bigl(-x_1-Rx_2+E\bigr),
  \label{eq:bi-diodo}
\end{equation}
con $x_1=v_C$ en V, $x_2=i_L$ en mA, $h$ la característica de la \cref{sec:un-diodo}, $E=1.2$\,V, $R=1.5$\,k$\Omega$, $C=2$\,pF, $L=5$\,$\mu$H y $t$ en ns.

Los equilibrios cumplen $x_2=h(x_1)$ y $x_2=(E-x_1)/R$: son los mismos tres cortes con la recta de carga del capítulo~2, $Q_1=(0.063,\,0.758)$, $Q_2=(0.285,\,0.610)$ y $Q_3=(0.884,\,0.210)$. La jacobiana es
\[
  A=\begin{bmatrix}-h'(x_1)/C&1/C\\-1/L&-R/L\end{bmatrix},
  \qquad
  \det A=\frac{1+Rh'(x_1)}{CL},
\]
de modo que la silla aparece exactamente donde $h'(x_1)<-1/R$, la misma condición del capítulo~2. Numéricamente:
\begin{center}\small
\begin{tabular}{lccc}
\toprule
 & $h'(\eq_1)$ & $\lambda_1,\lambda_2$ & tipo\\
\midrule
$Q_1$ & $7.24$ & $-3.59,\ -0.33$ & nodo estable\\
$Q_2$ & $-3.64$ & $1.77,\ -0.25$ & silla\\
$Q_3$ & $2.87$ & $-1.34,\ -0.40$ & nodo estable\\
\bottomrule
\end{tabular}
\end{center}
El circuito es biestable: es una celda de memoria. La inductancia no cambia los equilibrios ni su estabilidad, pero sí el \emph{camino}: en el modelo de primer orden el estado se mueve sobre la recta de $v_C$; aquí cada trayectoria es una curva en el plano, y cuál de los dos nodos se alcanza ya no depende solo de $v_C(0)$, sino también de $i_L(0)$.

\begin{figure}[ht]
\centering
\includegraphics[width=0.85\linewidth]{cap04/f_diodo}
\caption{Retrato del diodo túnel \eqref{eq:bi-diodo}. Las nulclinas son la característica $x_2=h(x_1)$ (coral) y la recta de carga (ámbar). En negro, la variedad estable de la silla $Q_2$: separa las condiciones iniciales que terminan en $Q_1$ de las que terminan en $Q_3$.}
\label{fig:bi-diodo}
\end{figure}

\codigorepo{cap04_bidimensionales/C4_Ej10_DiodoTunel.m}{diodo}{Equilibrios, jacobianas y retrato del diodo túnel de segundo orden}

\begin{pausa}
En el diodo túnel con $v_C(0)=0.4$\,V, a la derecha de $Q_2$, ¿termina siempre el circuito en $Q_3$, como en el modelo de primer orden?
\respuesta{No. Con $i_L(0)=0$ termina en $Q_1$: la corriente del inductor arranca en cero, el diodo descarga el condensador y $v_C$ cruza hacia la izquierda de $Q_2$ antes de que $i_L$ alcance a sostenerlo. Hace falta $i_L(0)\gtrsim0.189$\,mA para terminar en $Q_3$. Al revés, con $v_C(0)=0.15$\,V (a la izquierda de $Q_2$) e $i_L(0)=1.5$\,mA se termina en $Q_3$. La frontera entre las dos cuencas es la variedad estable de $Q_2$ (\cref{fig:bi-diodo}), y depende de los dos estados.}
\end{pausa}

%% ---------------------------------------------------------------------
\section{Variedades estable e inestable}
\label{sec:bi-variedades}
%% ---------------------------------------------------------------------

En una silla lineal, $E^s$ es la única recta de condiciones iniciales que llegan al origen. En una silla no lineal, ¿qué conjunto juega ese papel? No puede ser una recta en general, pero Hartman--Grobman sugiere que es la imagen deformada de una: una curva.

\begin{definicion}[label=def:bi-variedades]{Variedades estable e inestable}
Para un equilibrio $\eq$ de $\dot x=f(x)$, la \clave{variedad estable} y la \clave{inestable} son
\[
  W^s(\eq)=\bigl\{x_0:\ \flujo t{x_0}\to\eq\ \text{cuando}\ t\to+\infty\bigr\},\qquad
  W^u(\eq)=\bigl\{x_0:\ \flujo t{x_0}\to\eq\ \text{cuando}\ t\to-\infty\bigr\}.
\]
Las versiones \emph{locales} $W^s_{\mathrm{loc}}$ y $W^u_{\mathrm{loc}}$ exigen además que la solución se quede en una vecindad pequeña de $\eq$.
\end{definicion}

\begin{teorema}[label=thm:bi-variedadestable]{Variedad estable \citep[\S2.7]{Per01}}
Sea $f\in\Cont k$ ($k\ge1$) y $\eq$ un equilibrio hiperbólico. Entonces $W^s_{\mathrm{loc}}(\eq)$ y $W^u_{\mathrm{loc}}(\eq)$ son variedades de clase $\Cont k$, de dimensiones $\dim E^s$ y $\dim E^u$, \emph{tangentes} en $\eq$ a $E^s$ y $E^u$ (los subespacios de la linealización trasladados a $\eq$). Son invariantes, y $W^s$, $W^u$ se obtienen de las locales siguiendo el flujo hacia atrás y hacia adelante.
\end{teorema}

En el plano, para una silla, $W^s$ y $W^u$ son curvas que pasan por el equilibrio: las cuatro ramas que salen de él se llaman \clave{separatrices}.

\begin{ejemplo}[label=ej:bi-variedad]{Una variedad estable que se puede escribir}
Halle $W^s$ y $W^u$ del origen de $\dot x_1=-x_1$, $\dot x_2=x_2+x_1^2$.
\solucion La linealización es $\diag(-1,1)$: silla, con $E^s$ el eje $x_1$ y $E^u$ el eje $x_2$. El eje $x_2$ es invariante ($x_1=0\Rightarrow\dot x_1=0$) y sobre él $\dot x_2=x_2$: $W^u$ es el eje $x_2$. Para $W^s$ se busca una curva invariante $x_2=h(x_1)$ tangente al eje $x_1$. Invariancia significa que la derivada de $x_2-h(x_1)$ a lo largo de las soluciones se anula sobre la curva:
\[
  \dot x_2-h'(x_1)\dot x_1=h(x_1)+x_1^2+x_1h'(x_1)=0 .
\]
Con $h=cx_1^2$: $c+1+2c=0$, $c=-\tfrac13$. Entonces $W^s=\{x_2=-\tfrac13x_1^2\}$, una parábola tangente a $E^s$ en el origen. La solución explícita lo confirma: $x_1=x_1(0)e^{-t}$ y
\[
  x_2(t)=\Bigl(x_2(0)+\tfrac13x_1(0)^2\Bigr)e^{t}-\tfrac13x_1(0)^2e^{-2t},
\]
que tiende a $0$ si y solo si el paréntesis es cero.
\end{ejemplo}

\begin{figure}[ht]
\centering
\includegraphics[width=0.75\linewidth]{cap04/f_variedad}
\caption{El \cref{ej:bi-variedad}. $W^s$ (verde) es una parábola tangente al eje $x_1$ ($E^s$, gris punteado). La solución que arranca en $(1,0)\in E^s$ no llega al origen: sube y escapa.}
\label{fig:bi-variedad}
\end{figure}

\begin{trampa}[label=tr:bi-variedad]{``Las condiciones iniciales sobre $E^s$ llegan a la silla''}
\plausible{Para la linealización es cierto, y cerca del equilibrio el sistema se parece a la linealización.}
En el \cref{ej:bi-variedad}, $E^s$ es el eje $x_1$, y la solución desde $(1,0)$ tiene $x_2(t)=\tfrac13(e^t-e^{-2t})\to\infty$. La variedad estable es solo \emph{tangente} a $E^s$ en el equilibrio; lejos de él (y en este ejemplo, en cualquier punto distinto del origen) se separa de la recta.
\correcto{$W^s$ es una curva tangente a $E^s$ en $\eq$. Para usarla hay que calcularla: con una serie, como en el ejemplo, o numéricamente, integrando hacia atrás en el tiempo desde puntos muy cercanos a $\eq$ sobre $E^s$.}
\end{trampa}

\subsection{Las separatrices como fronteras}

La variedad estable de una silla no la puede cruzar ninguna trayectoria (estaría sobre ella), y las condiciones iniciales a un lado y al otro se separan cerca de la silla, siguiendo las dos ramas de $W^u$. Por eso \alerta{$W^s$ de una silla suele ser la frontera entre dos cuencas de atracción}.

En Duffing \eqref{eq:bi-duffing}, $W^u$ del origen tiene dos ramas: una cae en espiral a $(1,0)$ y la otra a $(-1,0)$. $W^s$ también tiene dos ramas, que vienen del infinito enrollándose: las condiciones iniciales entre ellas pertenecen a una u otra cuenca en franjas alternas (\cref{fig:bi-duffing}). Con energía grande la viga pasa varias veces sobre la barrera antes de quedar atrapada, y un cambio pequeño en la condición inicial cambia el pozo final. Las franjas son cada vez más delgadas lejos del origen, pero el sistema no es caótico: las fronteras son curvas suaves.

\begin{figure}[ht]
\centering
\includegraphics[width=\linewidth]{cap04/f_duffing}
\caption{Duffing con $\delta=0.5$. Izquierda: cuencas de $(1,0)$ (azul) y $(-1,0)$ (coral), separadas por $W^s$ del origen (negro); en verde, $W^u$. Derecha: $x_1(t)$ desde $(1.5,0)$ y desde $(2,0)$: la segunda tiene energía para pasar la barrera y termina en el otro pozo.}
\label{fig:bi-duffing}
\end{figure}

\video{V4_2_cuencas_duffing}{La variedad estable crece y dibuja las cuencas}
\codigorepo{cap04_bidimensionales/C4_Ej11_Variedades.m}{variedad}{$W^s$ y $W^u$ integrando hacia atrás y hacia adelante desde la silla}

\begin{pausa}
¿Por qué se integra \emph{hacia atrás} en el tiempo para dibujar $W^s$?
\respuesta{Porque hacia adelante $W^s$ es repulsora en la dirección transversal: un error numérico de $10^{-12}$ crece como $e^{\lambda_ut}$ y la solución se despega de la curva. Hacia atrás en el tiempo, $W^s$ se vuelve atractora en la dirección transversal ($W^s$ del sistema original es $W^u$ del sistema con el tiempo invertido) y los errores se amortiguan.}
\end{pausa}

%% ---------------------------------------------------------------------
\section{La variedad central}
\label{sec:bi-central}
%% ---------------------------------------------------------------------

Queda el caso que Hartman--Grobman deja abierto: la linealización tiene valores propios en el eje imaginario. Si además tiene alguno con parte real positiva, el equilibrio es inestable y no hay nada más que decidir (\cref{thm:ly-indirecto}). El caso difícil es entonces: \alerta{$k$ valores propios con parte real nula y los otros $m=n-k$ con parte real negativa}. La idea es que las direcciones estables se ``apagan'' exponencialmente rápido, y lo que decide la estabilidad pasa en una variedad de dimensión $k$, tangente a $E^c$, donde la dinámica es lenta.

Con un cambio de coordenadas lineal que lleve $A$ a la forma diagonal por bloques, el sistema queda
\begin{equation}
  \dot y=A_1y+g_1(y,z),\qquad \dot z=A_2z+g_2(y,z),\qquad y\in\R^k,\ z\in\R^m,
  \label{eq:bi-central}
\end{equation}
donde los valores propios de $A_1$ tienen parte real nula, $A_2$ es Hurwitz, y $g_1,g_2$ son de clase $\Cont2$ y se anulan con sus primeras derivadas en el origen.

\begin{teorema}[label=thm:bi-central]{Variedad central \citep[Teoremas~8.1 y~8.2]{Kha02}, \citep{Car81}}
\begin{enumerate}
  \item \emph{Existencia.} Existe una función $h$, definida cerca de $y=0$, con $h(0)=0$ y $\jac hy(0)=0$, tal que la superficie $z=h(y)$ es invariante para \eqref{eq:bi-central}: es una \clave{variedad central}.
  \item \emph{Reducción.} La dinámica sobre ella está dada por el \clave{sistema reducido}
  \begin{equation}
    \dot y=A_1y+g_1\bigl(y,h(y)\bigr) .
    \label{eq:bi-reducido}
  \end{equation}
  Si el origen de \eqref{eq:bi-reducido} es asintóticamente estable (estable, inestable), el origen de \eqref{eq:bi-central} también lo es.
\end{enumerate}
\end{teorema}

\notamargen{La prueba de la parte 2 usa una función de Lyapunov que combina la del sistema reducido con la parte estable $z-h(y)$; es la técnica del \cref{cap:lyapunov}. Ver \citet[\S8.1]{Kha02}.}

Para encontrar $h$ se usa la invariancia, igual que con la variedad estable: la derivada de $z-h(y)$ a lo largo de las soluciones debe anularse sobre $z=h(y)$, es decir,
\begin{equation}
  \jac hy(y)\Bigl[A_1y+g_1\bigl(y,h(y)\bigr)\Bigr]-A_2h(y)-g_2\bigl(y,h(y)\bigr)=0 .
  \label{eq:bi-pde}
\end{equation}
Casi nunca se puede resolver exactamente, pero basta aproximar: si una función $\phi$ con $\phi(0)=0$, $\phi'(0)=0$ satisface \eqref{eq:bi-pde} hasta términos de orden $\norm y^p$, entonces $h(y)=\phi(y)+O(\norm y^p)$ \citep[Teorema~8.3]{Kha02}. Con $y$ escalar se propone $h(y)=h_2y^2+h_3y^3+\cdots$ y se igualan coeficientes.

\begin{ejemplo}[label=ej:bi-central]{El ejemplo de siempre}
Determine la estabilidad del origen de $\dot y=yz$, $\dot z=-z+ay^2$, con $a\neq0$.
\solucion La linealización es $\diag(0,-1)$: un valor propio cero, que la linealización no decide. El sistema ya está en la forma \eqref{eq:bi-central} con $A_1=0$, $A_2=-1$, $g_1=yz$, $g_2=ay^2$. La ecuación \eqref{eq:bi-pde} es
\[
  h'(y)\,y\,h(y)+h(y)-ay^2=0 .
\]
Con $h=h_2y^2+h_4y^4+O(y^6)$ (solo potencias pares, por la simetría $y\to-y$): el coeficiente de $y^2$ da $h_2-a=0$; el de $y^4$, $2h_2^2+h_4=0$. Luego $h(y)=ay^2-2a^2y^4+O(y^6)$ y el sistema reducido es
\[
  \dot y=y\,h(y)=ay^3-2a^2y^5+O(y^7) .
\]
Si $a<0$, el origen del sistema reducido es asintóticamente estable (como $\dot y=-y^3$), y por el \cref{thm:bi-central} también el del sistema completo; si $a>0$, ambos son inestables. Basta el primer término; el segundo no cambia la conclusión.
\end{ejemplo}

\begin{figure}[ht]
\centering
\includegraphics[width=\linewidth]{cap04/f_central}
\caption{El \cref{ej:bi-central}. Las trayectorias (azul) caen rápido sobre la variedad central $z\approx ay^2$ (verde) y después se mueven despacio sobre ella: hacia el origen si $a=-1$, hacia afuera si $a=1$. La recta $z=0$ ($E^c$, gris) no es invariante.}
\label{fig:bi-central}
\end{figure}

\video{V4_3_variedad_central}{Caer rápido, moverse despacio}
\codigorepo{cap04_bidimensionales/C4_Ej12_VariedadCentral.m}{central}{Coeficientes de $h$ con la ecuación de invariancia y comparación con la simulación}

\begin{trampa}[label=tr:bi-zcero]{``$z$ es la variable rápida y estable: se reemplaza por $z=0$''}
\plausible{$z$ decae como $e^{-t}$, así que después de un transitorio $z\approx0$ y queda $\dot y\approx0$.}
Con $z=0$ el sistema reducido sería $\dot y=0$, para cualquier $a$, y no se podría distinguir el caso estable del inestable. $z$ no tiende a $0$: tiende a $h(y)=ay^2+\cdots$, que es pequeño pero no nulo, y ese término de segundo orden, multiplicado por $y$, es justo el que decide. $E^c$ (la recta $z=0$) no es invariante; la variedad central sí.
\correcto{Se reemplaza $z$ por $h(y)$, calculada con la ecuación de invariancia hasta el orden que decida la estabilidad del sistema reducido.}
\end{trampa}

\begin{pausa}
Con $a=0$ el sistema es $\dot y=yz$, $\dot z=-z$. ¿Cuál es la variedad central y qué concluye el \cref{thm:bi-central}? Compruébelo con la solución exacta.
\respuesta{$z=0$ es invariante y tangente a $E^c$: $h=0$, y el sistema reducido es $\dot y=0$. Su origen es estable pero no asintóticamente estable, y el teorema da lo mismo para el sistema completo. En efecto, $z=z_0e^{-t}$ y $y=y_0\exp\bigl(z_0(1-e^{-t})\bigr)\to y_0e^{z_0}$: las soluciones se quedan cerca, pero no vuelven al origen.}
\end{pausa}

\begin{trampa}[label=tr:bi-unica]{``La variedad central es única''}
\plausible{Se calcula resolviendo una ecuación, y la serie de $h$ queda determinada coeficiente por coeficiente.}
Para $\dot y=y^2$, $\dot z=-z$, las trayectorias cumplen $\dif z/\dif y=-z/y^2$, es decir, $z=c\,e^{1/y}$. Para $y<0$, $e^{1/y}\to0$ más rápido que cualquier potencia de $y$ cuando $y\to0^-$. Cada curva ``$z=ce^{1/y}$ si $y<0$, $z=0$ si $y\ge0$'' es invariante, tangente a $E^c$ y de clase $\Cont\infty$: hay infinitas variedades centrales, y todas tienen la misma serie de Taylor (cero) en el origen (\cref{fig:bi-unica}).
\correcto{La variedad central puede no ser única, pero todas coinciden hasta todos los órdenes en el origen, y el \cref{thm:bi-central} da la misma conclusión con cualquiera de ellas.}
\end{trampa}

\begin{figure}[ht]
\centering
\includegraphics[width=0.7\linewidth]{cap04/f_unica}
\caption{$\dot y=y^2$, $\dot z=-z$: para $y<0$, cada trayectoria $z=ce^{1/y}$ llega al origen tangente al eje $y$ y se puede continuar con $z=0$ para $y\ge0$. Todas son variedades centrales.}
\label{fig:bi-unica}
\end{figure}

%% ---------------------------------------------------------------------
\section{Para llevar}
\label{sec:bi-cierre}
%% ---------------------------------------------------------------------

\begin{parallevar}
\begin{enumerate}
  \item \textbf{Plano de fase:} el campo $f(x)$ en una malla, las trayectorias tangentes a él y las nulclinas ($f_1=0$: campo vertical; $f_2=0$: horizontal). Los equilibrios son los cruces de una nulclina de $x_1$ con una de $x_2$. El retrato no muestra el tiempo.
  \item \textbf{Lineales:} $x(t)=e^{At}x_0$ es la única solución. Con vectores propios, $x=Pz$ desacopla el sistema; cada estado mezcla todos los modos $e^{\lambda_it}$.
  \item \textbf{Clasificación en $\R^2$} ($\delta=\det A$, $\tau=\tr A$): silla si $\delta<0$; nodo si $\delta>0$ y $\tau^2\ge4\delta$; foco si $\tau^2<4\delta$ y $\tau\neq0$; centro si $\delta>0$ y $\tau=0$. El signo de $\tau$ da la estabilidad; el de $a_{21}$, el sentido de giro. Los nodos llegan por la dirección lenta.
  \item \textbf{Robustez:} nodos, focos y sillas sobreviven a perturbaciones pequeñas; centros, nodos degenerados y valores propios nulos, no.
  \item \textbf{Hartman--Grobman:} si $\eq$ es hiperbólico, el retrato no lineal cerca de $\eq$ es el lineal deformado por un homeomorfismo. Si no es hiperbólico, la linealización no decide.
  \item \textbf{Variedades:} $W^s$ y $W^u$ de una silla son curvas tangentes a $E^s$ y $E^u$; $W^s$ separa cuencas de atracción. Se calculan con la ecuación de invariancia o integrando hacia atrás.
  \item \textbf{Variedad central:} con valores propios nulos y los demás estables, la estabilidad se decide en el sistema reducido $\dot y=A_1y+g_1(y,h(y))$, con $h$ aproximada por la ecuación de invariancia.
\end{enumerate}
\end{parallevar}

\begin{trampa}[title={\faBug\hspace{0.45em}Las trampas del capítulo, en una caja}]{}
\begin{itemize}
  \item Un cruce de dos ramas de la misma nulclina no es un equilibrio (\cref{tr:bi-nulclinas}); las trayectorias no llegan a un equilibrio en tiempo finito (\cref{tr:bi-llega}).
  \item El signo de $a_{ii}$ no decide la estabilidad de $x_i$ (\cref{tr:bi-diagonal}).
  \item Los nodos se alcanzan por la dirección lenta (\cref{tr:bi-rapida}); el sentido de giro se lee en $a_{21}$, no en la forma de Jordan (\cref{tr:bi-giro}).
  \item Hartman--Grobman no distingue nodos de focos (\cref{tr:bi-hgfoco}); un centro lineal no es un centro no lineal (\cref{tr:bi-centro}).
  \item Duffing sin forzamiento no es caótico (\cref{tr:bi-caos}).
  \item $E^s$ no es $W^s$ (\cref{tr:bi-variedad}); $E^c$ no es la variedad central y $z=0$ no es $z=h(y)$ (\cref{tr:bi-zcero}); la variedad central puede no ser única (\cref{tr:bi-unica}).
\end{itemize}
\end{trampa}

\begin{modelo}[DiodoTunel2D.slx]
\cmd{crearModelo\_DiodoTunel2D.m} arma el circuito \eqref{eq:bi-diodo} con dos integradores y un bloque \texttt{Fcn} para la característica $h(v)$, con el tiempo en ns. Simula desde $(v_C,i_L)=(0.4,0)$ y $(0.4,0.3)$, que terminan en equilibrios distintos aunque empiezan con el mismo voltaje, y compara con \texttt{ode45}. Deje $E$, $R$, $L$, $C$ y $x_0$ en el workspace y busque la frontera entre las cuencas.
\end{modelo}

\begin{profundiza}[Correcciones respecto a la versión anterior (Handout3 y las diapositivas SisBidim)]
(i) La solución del ejemplo $\dot x_1=-x_1-3x_2$, $\dot x_2=2x_2$ es $x_1=x_1(0)e^{-t}+x_2(0)(e^{-t}-e^{2t})$, no $(e^{-t}-e^{-2t})$. (ii) En la linealización faltaba el factor $(x-\eq)$: $f(x)\approx Df(\eq)(x-\eq)$. (iii) En el caso de valores propios reales distintos, $J=P^{-1}AP$ (no $PAP^{-1}$) y $\dot z_2=\lambda_2z_2$. (iv) ``Cuando un valor propio es cero el origen es un equilibrio aislado'': es al revés, hay una recta de equilibrios. (v) En el foco: $\phi=\tan^{-1}(z_2/z_1)$, y ``si $\alpha=0$ la trayectoria tiende al origen'' es falso (es un centro). (vi) El ejemplo ``sin equilibrios'' con $\alpha>0$: hace falta $\alpha>1$. (vii) ``Duffing sin amortiguamiento tiene comportamiento caótico'': es integrable (\cref{tr:bi-caos}); con $\delta>0$, $(\pm1,0)$ son focos solo si $\delta<2\sqrt2$. La empresa de Duffing era la Ölwerke Stern-Sonneborn (Ossag), y fue Max Planck, no Plank. (viii) La prueba de la proposición de equilibrios aislados tenía la desigualdad al revés ($\norm{x-\eq}>\delta$) y mezclaba el caso variante en el tiempo; se reescribió. (ix) Hartman--Grobman: el homeomorfismo no es único, como decía la versión anterior, y no distingue nodos de focos. (x) En el ejemplo $\dot x_2=-x_1-\varepsilon x_1^2x_2$, con $\varepsilon>0$ el origen es asintóticamente estable (LaSalle), no solo estable. (xi) La prueba de la reducción a la variedad central no usa ``funciones tipo Kampke'' sino una función de Lyapunov, y el teorema necesita que no haya valores propios con parte real positiva. (xii) El teorema fundamental: se eliminó el paso $A^{-1}A$ (que suponía $A$ invertible) y el ``teorema de Moore''. (xiii) Se agregaron: nulclinas, el tiempo de llegada, el sentido de giro, la clasificación completa y su prueba, la robustez de los retratos, el diodo túnel de segundo orden, las variedades estable e inestable, las cuencas de Duffing, la no unicidad de la variedad central y el término de orden cuatro de $h$.
\end{profundiza}

%% ---------------------------------------------------------------------
\section{Problemas}
\label{sec:bi-problemas}
%% ---------------------------------------------------------------------

Los problemas marcados con \etiquetaM{} requieren \matlab/\simulink; los marcados con \etiquetaEstrella{} son de profundización. Cada uno tiene una trampa: búsquela antes de dar la respuesta.

\subsection*{Plano de fase y sistemas lineales}

\begin{problema}[conejos y ovejas]\label{prob:bi-competencia}
Dos especies compiten por el mismo pasto: $\dot x_1=x_1(3-x_1-2x_2)$, $\dot x_2=x_2(2-x_1-x_2)$, con $x_1,x_2\ge0$ \citep[\S6.4]{Str15}. (a) Dibuje las nulclinas. (b) Halle los cuatro equilibrios y clasifíquelos. (c) Existe un equilibrio de coexistencia: ¿coexisten las especies?
\end{problema}

\begin{problema}[la exponencial a mano]\label{prob:bi-exponencial}
(a) Calcule $e^{At}$ con la serie para $A=\bigl[\begin{smallmatrix}0&1\\0&0\end{smallmatrix}\bigr]$. (b) Use (a) para calcular $e^{Jt}$ con $J=\bigl[\begin{smallmatrix}\lambda&1\\0&\lambda\end{smallmatrix}\bigr]$. ¿Por qué se puede separar $e^{(\lambda I+N)t}=e^{\lambda t}e^{Nt}$? (c) Compare con la solución del nodo impropio.
\end{problema}

\begin{problema}[clasificar]\label{prob:bi-clasificar}
Clasifique el origen y, si gira, diga en qué sentido: (a) $\bigl[\begin{smallmatrix}1&2\\3&4\end{smallmatrix}\bigr]$; (b) $\bigl[\begin{smallmatrix}-2&1\\-1&-2\end{smallmatrix}\bigr]$; (c) $\bigl[\begin{smallmatrix}2&1\\-1&0\end{smallmatrix}\bigr]$; (d) $\bigl[\begin{smallmatrix}1&-2\\1&-1\end{smallmatrix}\bigr]$.
\end{problema}

\begin{problema}[por dónde llega]\label{prob:bi-lenta}
Para $A=\bigl[\begin{smallmatrix}-3&1\\0&-1\end{smallmatrix}\bigr]$: (a) halle los vectores propios; (b) ¿en qué dirección llegan al origen casi todas las trayectorias? (c) ¿Cuáles no? (d) \etiquetaM{} Verifíquelo con \cmd{retratoFase}.
\end{problema}

\subsection*{Linealización}

\begin{problema}[el péndulo, otra vez]\label{prob:bi-pendulo}
Para $\dot x_1=x_2$, $\dot x_2=-\sin x_1-bx_2$: (a) clasifique los equilibrios $(k\pi,0)$ para $b=0.5$; (b) ¿para qué $b>0$ es el origen un nodo? (c) Con $b=0$, ¿qué dice Hartman--Grobman del origen? ¿Qué dice la energía $E=1-\cos x_1+\tfrac12x_2^2$?
\end{problema}

\begin{problema}[un centro que no lo es]\label{prob:bi-hopf}
Para $\dot x_1=\mu x_1-x_2-x_1(x_1^2+x_2^2)$, $\dot x_2=x_1+\mu x_2-x_2(x_1^2+x_2^2)$: (a) clasifique el origen por linealización para $\mu<0$, $\mu=0$ y $\mu>0$; (b) pase a polares y estudie $\dot r$; (c) ¿qué es el origen cuando $\mu=0$? ¿Qué aparece cuando $\mu>0$? (Este es el anuncio de los \cref{cap:orbitas,cap:bifurcaciones}.)
\end{problema}

\begin{problema}[la variedad inestable]\label{prob:bi-inestable}
Para $\dot x_1=x_1$, $\dot x_2=-x_2+x_1^2$: (a) clasifique el origen; (b) halle $W^s$ y $W^u$; (c) ¿llega al origen hacia atrás en el tiempo la solución que empieza en $(1,0)$?
\end{problema}

\subsection*{Variedad central y simulación}

\begin{problema}[cuando el primer término se anula]\label{prob:bi-central}
Para $\dot y=ay^3+yz$, $\dot z=-z+y^2$: (a) halle $h(y)$ hasta orden cuatro; (b) escriba el sistema reducido; (c) clasifique el origen para $a<-1$, $a>-1$ y $a=-1$. En el último caso, ¿es suficiente el término cúbico?
\end{problema}

\begin{problema}[la frontera de las cuencas]\label{prob:bi-cuenca}
\etiquetaM{} Con \cmd{crearModelo\_DiodoTunel2D.m}: (a) simule desde $(0.4,i_L(0))$ para $i_L(0)=0,\ 0.1,\ 0.2,\ 0.3$\,mA; (b) encuentre con bisección el valor de $i_L(0)$ donde cambia el equilibrio final; (c) compare con la variedad estable de $Q_2$ que dibuja \cmd{C4\_Ej10\_DiodoTunel}.
\end{problema}

\begin{problema}[la dirección lenta, en general]\label{prob:bi-tangente}
\etiquetaEstrella{} Sea $A\in\R^{2\times2}$ con valores propios $\lambda_1<\lambda_2<0$ y vectores propios $v_1$, $v_2$. Muestre que si $x_0$ no está en la recta generada por $v_1$, entonces $x(t)/\norm{x(t)}\to\pm v_2/\norm{v_2}$. ¿Qué pasa con $x(t)/\norm{x(t)}$ en un foco?
\end{problema}

%% ---------------------------------------------------------------------
\section{Soluciones seleccionadas}
%% ---------------------------------------------------------------------

\paragraph{Problema~\ref{prob:bi-competencia}}
(a) Nulclinas de $x_1$: $x_1=0$ y $x_1+2x_2=3$; de $x_2$: $x_2=0$ y $x_1+x_2=2$. (b) $(0,0)$: $\diag(3,2)$, nodo inestable. $(0,2)$: $\bigl[\begin{smallmatrix}-1&0\\-2&-2\end{smallmatrix}\bigr]$, nodo estable. $(3,0)$: $\bigl[\begin{smallmatrix}-3&-6\\0&-1\end{smallmatrix}\bigr]$, nodo estable. $(1,1)$: $\bigl[\begin{smallmatrix}-1&-2\\-1&-1\end{smallmatrix}\bigr]$, $\delta=-1$: silla. (c) No: el equilibrio de coexistencia es una silla, y casi todas las condiciones iniciales terminan en $(3,0)$ o en $(0,2)$; la variedad estable de $(1,1)$ separa las dos cuencas. Que exista un equilibrio no quiere decir que se observe.

\paragraph{Problema~\ref{prob:bi-exponencial}}
(a) $A^2=0$, así que $e^{At}=I+At=\bigl[\begin{smallmatrix}1&t\\0&1\end{smallmatrix}\bigr]$. (b) $J=\lambda I+N$ con $N=A$, y $\lambda I$ conmuta con $N$: $e^{Jt}=e^{\lambda t}\bigl[\begin{smallmatrix}1&t\\0&1\end{smallmatrix}\bigr]$. (c) $z_1=(z_1(0)+z_2(0)t)e^{\lambda t}$, $z_2=z_2(0)e^{\lambda t}$: el factor $t$ es el que hace dar media vuelta a las trayectorias del nodo impropio.

\paragraph{Problema~\ref{prob:bi-clasificar}}
(a) $\delta=-2$: silla. (b) $\tau=-4$, $\delta=5$, $\tau^2-4\delta=-4$: foco estable, $\lambda=-2\pm j$; $a_{21}=-1$: horario. (c) $\tau=2$, $\delta=1$, $\tau^2=4\delta$: valor propio doble $1$ con un solo vector propio, $(-1,1)$: nodo impropio inestable (no es una estrella porque $A\neq I$). (d) $\tau=0$, $\delta=1$: centro, $\lambda=\pm j$; $a_{21}=1$: antihorario.

\paragraph{Problema~\ref{prob:bi-lenta}}
(a) $\lambda=-3$: $v_1=(1,0)$; $\lambda=-1$: $v_2=(1,2)$. (b) Tangentes a $(1,2)$, la dirección lenta. (c) Las que empiezan sobre el eje $x_1$ (la recta de $v_1$), que llegan por él.

\paragraph{Problema~\ref{prob:bi-pendulo}}
(a) $k$ par: $\lambda^2+0.5\lambda+1=0$, foco estable; $k$ impar: $\lambda^2+0.5\lambda-1=0$, silla. (b) $\tau=-b$, $\delta=1$: nodo si $b^2\ge4$, es decir, $b\ge2$. (c) Con $b=0$ la linealización es un centro y Hartman--Grobman no dice nada. Pero $\dot E=\sin x_1\,x_2+x_2(-\sin x_1)=0$: las trayectorias viven en curvas de nivel de $E$, que cerca del origen son curvas cerradas (el origen es un mínimo estricto de $E$). Es un centro no lineal, y lo decidió una cantidad conservada, no la linealización.

\paragraph{Problema~\ref{prob:bi-hopf}}
(a) Jacobiana $\bigl[\begin{smallmatrix}\mu&-1\\1&\mu\end{smallmatrix}\bigr]$: foco estable si $\mu<0$, centro si $\mu=0$, foco inestable si $\mu>0$. (b) $\dot r=\mu r-r^3$, $\dot\theta=1$. (c) Con $\mu=0$, $\dot r=-r^3$: el origen es asintóticamente estable aunque la linealización sea un centro. Con $\mu>0$, $\dot r=0$ en $r=\sqrt\mu$: aparece un ciclo límite estable (bifurcación de Hopf).

\paragraph{Problema~\ref{prob:bi-inestable}}
(a) $\diag(1,-1)$: silla, $E^u$ el eje $x_1$, $E^s$ el eje $x_2$. (b) El eje $x_2$ es invariante y allí $\dot x_2=-x_2$: $W^s$ es el eje $x_2$. Para $W^u$, $x_2=cx_1^2$ con $2c=-c+1$: $c=\tfrac13$, $W^u=\{x_2=\tfrac13x_1^2\}$. (c) No: $(1,0)$ está en $E^u$ pero no en $W^u$. La solución es $x_2(t)=\tfrac13(e^{2t}-e^{-t})$, que hacia atrás se va a $-\infty$.

\paragraph{Problema~\ref{prob:bi-central}}
(a) La ecuación es $h'(y)(ay^3+yh)+h-y^2=0$. Con $h=y^2+h_4y^4$: $2y(ay^3+y^3)+h_4y^4=0$, $h_4=-2(a+1)$. (b) $\dot y=(a+1)y^3-2(a+1)y^5+\cdots$. (c) $a<-1$: asintóticamente estable; $a>-1$: inestable. Con $a=-1$ el término cúbico se anula, y no basta: hay que mirar más órdenes. Aquí todos se anulan, porque $h=y^2$ es \emph{exacta} (compruébelo) y el sistema reducido es $\dot y=-y^3+y\cdot y^2=0$: el origen es estable pero no asintóticamente estable.

\paragraph{Problema~\ref{prob:bi-cuenca}}
Desde $(0.4,0)$ y $(0.4,0.1)$ el circuito termina en $Q_1$; desde $(0.4,0.2)$ y $(0.4,0.3)$, en $Q_3$. La bisección da $i_L(0)\approx0.189$\,mA, el punto donde la recta $v_C=0.4$ corta la variedad estable de $Q_2$.

\paragraph{Problema~\ref{prob:bi-tangente}}
Escriba $x_0=c_1v_1+c_2v_2$ con $c_2\neq0$. Entonces $x(t)=e^{\lambda_2t}\bigl(c_1e^{(\lambda_1-\lambda_2)t}v_1+c_2v_2\bigr)$ y, como $\lambda_1-\lambda_2<0$, el paréntesis tiende a $c_2v_2$: la dirección tiende a $\sgn(c_2)\,v_2/\norm{v_2}$. En un foco, $x(t)/\norm{x(t)}$ gira sin parar y no tiene límite.

\respuestaspausas

